\scanpage{374}
\mspage{290}
% archive sheet number: 373

Alors \unsure{pour} $\pi\geq\pi'_1$ tel que $\pi^{\pi'_0}=\pi$,
$\underline S$ est $\pi$-accessible dans $M$ (4.13.2. b et c).

Supposons p. ex.
\[
\pi'_0\geq\pi_0
\]
(pratiquement ce sera tjrs vérifié) et de plus
\[
\pi'_0\geq\deleted{\unsure{\pi_1=\operatorname{Sup}}}\operatorname{Card}\operatorname{tot}M_{\pi_0}
\]
i.e. supposons
\[
\pi'_0\geq\operatorname{Sup}(\pi_0,\operatorname{Card}\operatorname{tot} M_{\pi_0})\equiv\pi_1.
\]
Alors \deleted{\unsure{comme} $\unsure{\in S_{\pi'_0}}$}
\[
\deleted{\pi'_1=\operatorname{Sup}(\pi'_0,\operatorname{Card}\unsure{\operatorname{Fct}|\operatorname{tot}|\operatorname{Fl}}S_{\pi'_0})}\ \leftarrow
\]
\[
\pi'_1=\operatorname{Sup}(\pi'_0,\operatorname{card}\operatorname{tot}S_{\pi'_0})\leq\operatorname{Sup}(\pi'_0,\operatorname{card}\operatorname{tot}M_{\pi'_0})\leq 2^{({\pi'_0}^{\pi_0})}.
\]
(4.7.7.1 p. \unsure{173|177}). Donc on trouve ceci :

\uline{Corollaire} 4.13.6. Soit $\pi'_0$ un cardinal $\geq\pi_1$,
\inserted{$\operatorname{Sup}(\pi_0,\operatorname{card}\operatorname{tot} M_{\pi_0})$}
et soit $\underline S$ une sous-catégorie $\pi'_0$-accessible
dans $M$. Alors pour \deleted{tt} cardinal $\pi$
% slanted
\margin{\uline{NB} $\pi_0$ est le plus petit cardinal utile pour $M$}
tel que $\pi\geq 2^{({\pi'_0}^{\pi_0})}$, \inserted{et} tel que l'on
ait $\pi^{\pi'_0}=\pi$ (i.e. $\pi$ de la forme
$c^{\pi'_0}$, où $c$ un cardinal $\geq 2^{({\pi'_0}^{\pi_0})}$),
$\underline S$ est $\pi$-accessible dans $M$ ! (De plus,
ces cardinaux sont réguliers pour $\underline S$ \uuline{et}
pour $M$, donc $\underline S\approx\operatorname{Ind}_{\pi}(\underline S_{\pi})$, $M\cong\operatorname{Ind}_{\pi}(M_\pi)$.)

\uline{Enchaîner vers \unsure{nos|mes|n°s} 4.13.23. \unsure{etc.}}

% slanted
\margin{Mettre\unsure{!} ici le lemme \unsure{4.16.7.11|4.16.7.1)|4.16.7.II} signalé en marge p. 292}
